%%%%%%%%%%%%%%%%%%%%%%%%%%%%%%%%%%%%%%%%%%%%%%%%%%%%%%%%%%%%%%%%%%%%%%%%%%%%%%%%%%
%% This project aims to create a test template or exercise list from the        %%     
%% Federal University of Ceará (UFC).                                           %%
%% author: Maurício Moreira Neto - Doctoral student in Computer Science         %%
%% contacts:                                                                    %%
%%    e-mail: maumneto@ufc.br                                                   %%
%%    linktree:  https://linktr.ee/maumneto                                     %%
%%%%%%%%%%%%%%%%%%%%%%%%%%%%%%%%%%%%%%%%%%%%%%%%%%%%%%%%%%%%%%%%%%%%%%%%%%%%%%%%%%
\documentclass{ufcdocument}

\usepackage[utf8]{inputenc}
\usepackage[portuguese]{babel}
\usepackage{amssymb}

\begin{document}

       %% Informations that will be insert in the table header 
\def\course{Lógica para Ciência da Computação}
\def\prof{Esdras Lins Bispo Jr.}
\def\semester{2026.1}
\def\codeCourse{ICE0239}
\def\registration{}
\def\student{}
\def\graduate{Bacharelado em Ciência da Computação}
\def\theme{Avaliação da Competência LCC-C1.1(a):\\ \texttt{Conversão da linguagem simbólica para a natural.}
}
    %% Table with the header
    \makeheader
    
    %% Space for the instructions
    % \fbox{
    %     \parbox{\textwidth}{
    %         \begin{minipage}{\textwidth}
    %             \makeinstructions
    %             {
    %                 \begin{instlist}
    %                     \item A avaliação é individual e não é pesquisada.
    %                     \item Preencha o cabeçalho da folha pergunta com seus dados.
    %                     \item  Todas as folhas respostas devem conter o nome a a matrícula do aluno.
    %                     \item O preenchimento das respostas deve ser feito utilizando caneta (preta ou azul).
    %                 \end{instlist}
    %             }
    %         \end{minipage}
    %     }
    % }
    %% Space between the instructions and the questions.
    \vspace{0.5cm}
    
    \begin{question}
        \item Sejam as proposições $x$: ``A estudante cursa programação introdutória'', e $z$: ``A estudante consegue programar utilizando funções''. Traduza  a proposição $\sim (z$ $\leftrightarrow$ $\sim x)$ para a linguagem natural.
    \end{question}
 
    \vfill

    \begin{center}
    {\small
    \setlength{\tabcolsep}{10pt}
    \renewcommand{\arraystretch}{2}
    \begin{tabular}{|p{0.47\textwidth}|c|c|}
    \hline
    \multicolumn{3}{|c|}{\textbf{Rubrica --- C1.1: Conversão da linguagem simbólica para a natural}} \\
    \hline
    \textbf{Critérios} & \textbf{Em formação} & \textbf{Proficiente} \\
    \hline
    Compreensão do Enunciado & $\square$ & $\square$ \\
    \hline
    Leitura da Linguagem Simbólica & $\square$ & $\square$ \\
    \hline
    Escrita da Linguagem Natural & $\square$ & $\square$ \\
    \hline
    \end{tabular}
    }
    \end{center}

\end{document}